% ============================================================
%  Appendix A — Implementation and Code
% ============================================================
\chapter{Implementation and Code}
\label{app:code}

This appendix documents the software implementation of the analysis. The full
pipeline was written in Python~3.13 and organized as a sequence of scripts, each
corresponding to one stage of the workflow in Chapter~\ref{ch:methodology}.

\section{Software Environment}
\label{app:env}

The principal libraries used were \texttt{earthengine-api} and \texttt{geemap}
for data collection; \texttt{rasterio} and \texttt{geopandas} for geospatial
preprocessing; \texttt{numpy} and \texttt{pandas} for array and table handling;
and \texttt{scikit-learn}, \texttt{xgboost}, and \texttt{lightgbm} for model
training and evaluation. Figures were produced with \texttt{matplotlib} and
\texttt{contextily}.

\section{Pipeline Scripts}
\label{app:scripts}

Table~\ref{tab:scripts} lists the scripts in execution order and the role of each.

\begin{table}[H]
\centering
\caption{Pipeline scripts in execution order.}
\label{tab:scripts}
\small
\begin{tabular}{p{4.4cm} p{8.6cm}}
\toprule
Script & Function \\
\midrule
\texttt{collect\_data.py}        & Assemble the nine features in Google Earth Engine and export the multi-band GeoTIFF. \\
\texttt{preprocess.py}           & Mask to district boundaries, min--max normalize each band, quality-check. \\
\texttt{step1\_generate\_labels.py} & Weighted-overlay scoring and percentile thresholding into three classes. \\
\texttt{step2\_validate.py}      & Exploratory analysis: correlations, distributions, class balance. \\
\texttt{step3\_train\_models.py} & Train and evaluate RF, XGBoost, LightGBM, and the voting ensemble; 5-fold CV. \\
\texttt{step4\_predict\_map.py}  & Tiled full-plateau prediction; classified and probability rasters. \\
\texttt{step9\_spatial\_realism\_validation.py} & The four spatial-realism checks and summary figure. \\
\bottomrule
\end{tabular}
\end{table}

\section{Representative Code}
\label{app:snippets}

The weighted-overlay score that produced the training labels
(Equation~\ref{eq:overlay}) was implemented as a direct weighted sum of the
normalized bands, with negatively-related features inverted before weighting:

\begin{verbatim}
# Weighted overlay: composite groundwater potential score
weights = {
    'soil': 0.25, 'twi': 0.20, 'rainfall': 0.20, 'ndvi': 0.10,
    'slope': 0.10, 'ndwi': 0.05, 'elevation': 0.05, 'lst': 0.03,
    'aspect': 0.02,
}
invert = {'slope', 'elevation', 'lst'}   # higher value -> lower potential

score = np.zeros(n_pixels, dtype=np.float32)
for name, w in weights.items():
    f = bands[name]                      # normalized to [0, 1]
    score += w * (1.0 - f if name in invert else f)

# Percentile thresholds: 40% Low, 40% Medium, 20% High
low_t, high_t = np.percentile(score, [40, 80])
labels = np.where(score < low_t, 0, np.where(score < high_t, 1, 2))
\end{verbatim}

The XGBoost classifier was configured with the hyperparameters reported in
Section~\ref{sec:ml_models}:

\begin{verbatim}
from xgboost import XGBClassifier

model = XGBClassifier(
    n_estimators=300,
    learning_rate=0.1,
    max_depth=6,
    subsample=0.8,
    colsample_bytree=0.8,
    objective='multi:softprob',
    num_class=3,
    n_jobs=-1,
)
model.fit(X_train, y_train)
\end{verbatim}

Full-plateau prediction was carried out over $2048 \times 2048$ pixel tiles to
keep memory use bounded, with each tile classified and written back to the output
raster before the next was read.
